% 本文件由 LaTeX 排版 Agent 根据有效 translation.json 自动生成。
% author=Theodotus; work_id=0802; title=狄奥多图斯摘录

\chapter{狄奥多图斯摘录}

% structure: p0001 source=h1 target=omitted reason=page-title-already-rendered-by-parent

\\Emph{选自先知书}\par

火窑中的沙得辣客、米沙、阿贝得乃哥三人赞美天主时如此说：\\BibleQuote{诸天，请赞美上主，歌颂称扬他，直到永远！}然后：\\BibleQuote{上主的天使，请赞美上主；}然后：\\BibleQuote{天上所有的水，请赞美上主。}因此圣经将诸天和水归入纯洁的\\OriginalTerm{权能}{power}之列，正如创世记所示。既然\\Quote{权能}一词有多种含义，达尼尔接着又说：\\BibleQuote{愿所有的权能都赞美上主；}然后又说：\\BibleQuote{太阳和月亮，请赞美上主；}以及\\BibleQuote{天上的星辰，请赞美上主。凡敬拜\\Emph{他}的，请赞美上主，歌颂并承认万神之神，因为他的仁慈永远长存。}这是达尼尔书中所记，当三少年在火窑中赞美时的话。\par

\\BibleQuote{你坐在革鲁宾之上，俯视深渊，是应当赞美的。}达尼尔如此说，与哈诺客的话相符，哈诺客曾说：\\BibleQuote{并且我看见了各种物质。}因为\\OriginalTerm{深渊}{abyss}本质上是无限的，却被天主的权能所约束。这些物质本质，从中产生各属类及其物种，被称为深渊；因为你不应仅称水为深渊，尽管物质被喻为水，即深渊。\par

\\BibleQuote{在起初天主创造了天地，} \\BibleRef{创 1:1} 包括地上的和天上的事物。这真理主曾对欧瑟亚说：\\BibleQuote{你去娶一个妓女为妻，生妓女的子女；因为这地行淫，\\Emph{偏离}了上主。} \\BibleRef{欧 1:2} 这里他所说的不是元素——\\Emph{地}，而是居住在这元素中的人，即那些心存世俗的人。\par

至于圣子乃是开始\\Emph{或\\OriginalTerm{头}{head}}，欧瑟亚明确教导：\\BibleQuote{将来，在从前对他们说\\Quote{你们不是我的人民}的地方，他们要被称为永生天主的子女；犹大子民和以色列子民要聚集在同一处，立一个头在他们之上，并从地上上来，因为依次勒耳的日子是伟大的。} \\BibleRef{欧 1:10} 因为人信谁，就拣选谁。但人所信的是圣子，他就是头；因此他又补充说：\\BibleQuote{但是我要怜悯犹大的儿子们，并借他们的天主上主拯救他们。} \\BibleRef{欧 1:7} 然而拯救他们的救主就是天主子。所以他便是头。\par

圣神借欧瑟亚说：\\BibleQuote{我是你们的\\OriginalTerm{教师}{Instructor}；} \\BibleRef{欧 5:2} \\BibleQuote{你要在上主的山上吹号角，在高处扬声。} \\BibleRef{欧 5:8} 而\\OriginalTerm{圣洗}{baptism}本身，就是重生的记号，难道不是藉着救主的教导，逃离物质，一条巨大汹涌的河流，永远奔流并携带着我们吗？主因此引领我们脱离混乱，光照我们，将我们带入无影之光，这光已不再是物质的。\par

这条\\OriginalTerm{物质}{matter}的\\OriginalTerm{河流}{river}与\\OriginalTerm{海}{sea}被两位先知藉上主的权能劈开并分开，物质被约束，经由水的两次分开。两位著名的领袖，因着他们所相信的奇迹，遵行了天主的旨意，使义人得以脱离物质，先穿越它而行。其中一位统帅还被赋予了我们的救主之名。\par

\\OriginalTerm{重生}{regeneration}是藉水和圣神而成，正如整个受造界：\\BibleQuote{天主的圣神运行在\\OriginalTerm{深渊}{abyss}之上。} \\BibleRef{创 1:2} 为此缘故，救主虽本身无需领洗，却接受了\\OriginalTerm{圣洗}{baptism}，以便为那些将要重生的人祝圣所有的水。因此我们所洁净的不仅是身体，也是灵魂。这标志着我们无形部分的圣化，并从新的属灵受造中滤除那些混杂在灵魂中的不洁之灵。\par

\\OriginalTerm{天上的水}{water above the heaven}。既然圣洗是由水和圣神施行，作为一种保护对抗双重之\\OriginalTerm{火}{fire}——那抓住可见之物的火和那抓住不可见之物的火；并且必然地，存在一种非物质的水元素和一种物质的水元素，它便是对抗双重火的保护。地上的水洁净身体；而天上的水，因其非物质和无形，是圣神的象征，圣神是无形之物的净洗者，如同圣神之水，另一者则是身体之水。\par

天主出于\\OriginalTerm{良善}{goodness}，将\\OriginalTerm{敬畏}{fear}与良善相混合。因为他供给每个人有益的，如同医生对待病人，父亲对待悖逆的儿子：\\BibleQuote{不肯用杖打儿子的，是恨恶他。} \\BibleRef{箴 13:24} 主和他的宗徒们在敬畏与劳苦中行走。因此，当苦难降临在义人身上时，要么是主为先前所犯的罪而责备他，要么是为将来而保护他，要么是藉其权能不阻止外来的攻击——为了他和近人的益处，作为榜样。\par

那些居住在\\OriginalTerm{腐朽的身体}{corrupt body}中的人，如同乘坐一艘旧船航行的人，并非仰躺，而是不断祈祷，向天主伸开双手。\par

古人们极为痛苦，除非他们身体常有某种苦难。因为他们害怕，若今世不承受因无知而伴随肉身之众罪的惩罚，来世便要一并承受刑罚。因此他们宁可在此接受治疗。所以可怕的不是外在的疾病，而是罪；疾病\\Emph{由此而来}，且是灵魂之病，而非身体之病：\\BibleQuote{凡有\\OriginalTerm{血肉}{flesh}的尽都如\\OriginalTerm{草}{grass}，} \\BibleRef{依 40:6} 身体和外在的美好是暂时的；\\BibleQuote{但那看不见的才是永远的。} \\BibleRef{格后 4:18}\par

第十二。论知识，其中有些我们已经拥有；另一些，则借由我们已拥有的，我们坚定地希望\Emph{获得}。因为我们既未完全获得，也未完全缺乏。但我们仿佛已领受了永恒福乐和祖业的神印。为主预备道路的，是主的真福。因为祂说：\Quote{你们要寻求，}且要切切寻求，\Quote{天主的国，这一切自会加给你们；因为你们的父知道你们需要什么。}如此，祂不仅限制了我们的劳作，也限制了我们的挂虑。因为祂说：\Quote{你们不能因思虑而使身量增加一分一毫。}因为天主深知什么对我们有益、什么当缺少。因此，祂愿意我们弃绝世俗的挂虑，而充满那导向天主的。\Quote{因为我们叹息，渴望在脱去这必朽坏之前，穿上那不朽坏的。}因为当信德广传时，不信便不知所措。同样，知识和义德也是如此。因此，我们不仅要倒空灵魂，还要以天主充满它。因为其中不再有恶，因恶已止息；也未有善，因尚未领受善。而那既非善亦非恶的，就是虚无。\Quote{因为那打扫干净、空着的房屋，} \BibleRef{玛 12:44} 若没有放任何救恩的福分，那先前住在里面的不洁之神便会回来，并带上七个更恶的鬼。因此，在倒空灵魂的恶之后，我们必须以善神天主充满这蒙祂拣选的居所。因为当空室被充满后，便随之盖上印，使圣所得以为主守护。\par

第十三。\Quote{凭两三个见证人，一切事都得以确定。} \BibleRef{申 17:6} 即凭圣父、圣子、圣神，藉他们的见证和助佑，当遵守所吩咐的诫命。\par

第十四。禁食，按此词的意义，是指戒除食物。但食物并不能使我们更义或不义。然而，从奥秘意义上讲，它表明：正如生命藉食物维持，而缺乏食物是死亡的记号；同样，我们也当禁绝世俗之物，好死于世界，然后藉领受神圣的养料而为天主而活。尤其，禁食使灵魂摆脱物质，并使之与身体一同洁净轻盈，适于神圣之言。世俗的食物，便是先前的生活和诸罪；而神圣的食物，则是信德、望德、爱德、忍耐、知识、平安、节制。因为\Quote{饥渴慕义的人是有福的，因为他们要得饱饫。} \BibleRef{玛 5:6} 切盼这义的是灵魂，而非身体。\par

第十五。救主向信德的宗徒们表明，在一个他们不能洁净的附魔者身上，祈祷比信德更强有力，祂说：这一类事情是靠祈祷成就的。那信的人已从主获得罪赦；但那达到知识的人，因不再犯罪，便从自身获得其余之罪的赦免。\par

第十六。正如借着人施行医治、预言和神迹，是天主在其中工作；同样，诺斯替教导也是如此。因为天主藉人显示其权能。预言正确地宣告说：\Quote{我要给他们派遣一位拯救者。} \BibleRef{依 19:20} 因此，祂时而派遣先知，时而派遣宗徒，作人的拯救者。这样，天主藉着人行善。因为并非天主能做某些事、不能做另一些事；祂在任何事上从不无力。也不是有些事合其意愿、有些事违其意愿；有些事由祂而行，有些事由别人而行。相反，祂甚至借由我们而使我们存在，并借由我们而训练我们。\par

第十七。天主造了我们，我们先前并不存在。因为若我们先前存在，必当知道我们曾在何处，以及如何、为何来到此处。但若我们并非先前存在，那么天主就只是我们受造的唯一原因。因此，正如祂造了我们这不存在者；同样，如今我们既已成为受造，祂便因祂的恩宠拯救我们，只要我们显出自己堪当并领受。否则，祂便让我们归于适当的结局。因为祂是生者和死者的主。\par

第十八。但请看天主的能力，不仅是在人身上，从无中引出存在，并使已存在者按生命时程成长；而且祂也拯救那些相信的人，以适合各人的方式。如今祂也改变时辰、时期、果实和元素。因为这是唯一的天主，祂按各人衡量了事件的始和终。\par

第十九。从信德和敬畏迈向知识的人，知道如何称呼\Quote{主，主}，但不再如奴隶一般，他已学会说\Quote{我们的父}。他释放了产生恐惧的奴役之灵，藉爱而迈向义子之身，如今因爱敬畏祂从前所惧怕的。因为他不再出于恐惧而禁戒所当禁戒的，而是出于爱持守诫命。经上说：\Quote{圣神亲自与我们的心神一同作证，我们呼号：\InnerQuote{阿爸，父啊！}}\par

第二十。如今主以祂的宝血救赎了我们，使我们脱离昔日残忍的霸主，即我们的诸罪，正因这些罪，属灵的邪恶\Emph{权势}曾统治我们。因此，祂领我们进入父的自主——同为嗣子和朋友的儿女。\Quote{因为}主说：\Quote{凡遵行我父旨意的人，就是我的弟兄和同继承者。} \BibleRef{玛 12:50} \Quote{所以，不要称地上的人为你们的父。} \BibleRef{玛 23:9} 因为那是地上的主人。但在天上的是父，天上地下的一切家族都由祂得名。\BibleRef{弗 3:15} 因为爱统治甘愿的\Emph{心}，而恐惧统治不愿的。有一种恐惧是卑劣的；另一种，却如启蒙师引我们向善，带我们到基督面前，且是得救的。\par

XXI. 若有人对天主持有一个概念，这概念绝不相称于祂的尊贵。因为天主的尊贵能是什么呢？但他当尽可能设想一个伟大、不可理解且极美丽的光；不可接近的，涵括一切美善的能力，一切恩雅的德行；照顾万物，慈悲，无情感，良善；知晓一切预知一切，纯洁，甘饴，光明，无玷。\par

XXII. 既然灵魂的运动是自发的，天主的恩宠要求于它的，是它所拥有的，即意愿作为它对救恩的贡献。因为灵魂愿意自为善；\Emph{然而}，主却将此赐予它。因为它并非毫无知觉，以致像身体那样被带动。拥有是取的结果，取是意愿和欲求的结果；而持守所得，则是运用谨慎和能力的结果。因此，天主赋予了灵魂自由选择，为要指示其当尽之责，并使之在选择后，得以领受并持有。\par

XXIII. 正如主曾透过身体说话并医治，同样，先前藉着先知，如今藉着宗徒和教师。因为教会是主能力的服务者。因此祂于是取了人性，好藉此服务圣父的旨意。而且常时，那爱人者天主使自己穿上人，为人的救恩——先前在先知身上，如今在教会身上。因为同类服务同类，以达成同类的救恩，是合宜的。\par

XXIV. 因为我们是出于土的……凯撒是暂时的君王，其地上的肖像是那旧人，他已回归于旧人。那么，我们要将属于地上的事物归给他，就是我们按那属于地上的形象所背负的；并将属于天主的事物归给天主。因为每一种情欲在我们身上就如一封信、一个印记和一个记号。如今主用另一个印记、别的名字和字母来标记我们，即以信德代替不信，等等。这样，我们就从物质的被转化到属灵的，\BibleQuote{既然背负了那属于天上的形象。} \BibleRef{格前 15:49}\par

XXV. 若翰说：\BibleQuote{我固然用水洗你们，但在我以后要来的那一位，要用圣神和火洗你们。} \BibleRef{玛 3:11} 但祂并未用火洗任何人。但有些人，如赫拉克利所说，用火印在那些被封印者的耳朵上；如此理解宗徒的话：\BibleQuote{祂的簸箕在祂手中，要扬净祂的禾场；祂将麦子收入仓里，而将糠秕用不灭的火焚烧。} \BibleRef{玛 3:12} 因此，\Quote{用火}这一表达与\Quote{用圣神}相连；因为祂藉着圣神将麦子与糠秕分开，即从物质的壳分开；糠秕被簸开，被风吹走：同样，圣神也拥有分离物质力量的能力。既然有些事物是从非受造和不可毁灭的（即生命的种子）产生，麦子也被储存，而物质部分，只要它与优越部分结合，就存留；当它与优越部分分离时，就被毁灭；因为它存在于另一事物中。那么，这分离的元素就是圣神，毁灭的元素就是火：并且要理解为物质之火。但既然得救的像麦子，而在灵魂中生长的像糠秕，前者是无形的，后者是被分离的物质的；祂用圣神——一种稀薄、纯净，几乎比心智更精微的元素——来对抗无形的；而用火来对抗物质的，不是作为恶或坏的，而是作为强有力且能清除邪恶的。因为火被视为一种良善而强大的力量，能毁灭低贱者并保存更优者。因此，这火被先知们称为智慧之火。\par

XXVI. 同样，当日主被称为\Quote{吞噬之火}时，是因为要选取一个名称和记号，并非表示邪恶，而是表示能力。因为正如火是元素中最强大的，并能支配一切；同样，天主是全能和全能的，祂能持守、创造、制作、滋养、生长、拯救，对灵魂和身体拥有权能。既然火比所有元素更优越，同样，全能的主宰也高于众神、诸能力和一切宰制者。火的能力是双重的：一种能力有助于果实和动物的生长和成熟，太阳是其形象；另一种则用于消耗和毁灭，正如地上的火。因此，当日主被称为吞噬之火时，\Emph{祂被称为}一种强大而不可抗拒的能力，没有什么是不可能的，并且能够毁灭。\par

关于这样的能力，救主也说：\BibleQuote{我来把火投在地上，} \BibleRef{路 12:49} 指明一种能力，能净化圣物，但如他们所说，毁灭物质的东西；而按我们所说，则是惩戒性的。恐惧属于火，扩散属于光。\par

XXVII. 古人并不写作，因为他们既不愿因写作而占用本该用于专注他们所传授的内容的时间，也不愿将考虑应说什么的时间花在写作上。或许，他们确信写作的职能与教导的领域并不属于同一心智类型，便将此事让给有自然天赋的人。因为对于演说者来说，言语之流不受阻碍、奔腾不息，你可以仓促地捕捉它。但被读者反复检验的东西，经受严格的审查，被认为值得付出最大的努力，可以说是\Emph{口头}教导的书面确认，以及通过书面作品传递到后世的言语。因为受托于长者的东西，以文字言说，借助作者的帮助将其自身传递给那些将要阅读的人。因此，正如磁石排斥其他物质，只因亲和力而吸引铁；同样，书籍虽被许多人阅读，却只吸引那些能够理解它们的人。因为真理的言语，对一些人来说是\BibleQuote{愚蠢}（\BibleRef{格前 1:18}），对另一些人是\BibleQuote{绊脚石}（\BibleRef{格前 1:18}），但对少数人是\BibleQuote{智慧}（\BibleRef{格前 1:18}）。天主的权能也是如此。愿诺斯替派远离嫉妒。因为正是为此，他也发问：给不配的人，还是不给配得的人，哪一样更糟？出于他那满溢的传授之爱，他不仅冒险给予每个合格之人，有时也给予一个不配的人，若是他强求；并非因为他的恳求（因为他不好声誉），而是因为恳求者以丰富的恳求将心志转向信德。\par

XXVIII. 有些自称为诺斯替派的人，对自己家里的人比对外人更嫉妒。正如海洋向所有人敞开，但一人游泳，一人航行，一人捕鱼；陆地是公共的，但一人行走，一人耕种，一人打猎——有人探矿，有人造屋：同样，当圣经被阅读时，一人得助于信德，另一人得助于道德，第三人则因对事物的知识而摆脱迷信。那知道奥林匹克运动场的运动员，为训练而脱衣，竞赛，并成为胜利者，绊倒那些反对他科学方法的对手，奋战到底。因为科学知识对于灵魂的训练和行为的庄重都是必要的，它使信徒更积极、更敏锐地观察事物。因为没有基本教导就没有信德，同样，没有科学就没有理解。\par

XXIX. 对于救恩有用且必要的东西，例如\Emph{对圣父、圣子、圣神的知识}，以及对我们自己灵魂的知识，都是完全必需的；获得它们的科学说明既有益又必要。对于带头行善的人，丰富的经验是有益的，这样，别人所必须并博学地知道的事物中，没有一件能逃过他们的注意。对异端教导的解释，为探究的灵魂提供了另一种操练，并使门徒因事先在各种战争乐器上练习过而免于被诱惑偏离真理。\par

XXX. 诺斯替派的规范生活（正如他们说\Place{克里特}没有致命的动物）纯洁无瑕，远离一切恶行、恶念和恶言；不仅不恨任何人，而且超越嫉妒和仇恨，以及一切恶言和诽谤。\par

XXXI. 论长寿，人并非因活得久而被视为有福，而是因他在活着时也得以配得永生。他没有使人痛苦，只是用言语教导那心灵受伤的人，如同用有益的蜂蜜，既甜又刺。所以，诺斯替派尤其保持了与理性相一致的庄重。因为当激情被从整个灵魂中切除并剥除后，他从此与那最高贵者交往并生活，这高贵者现已变得纯洁，并得到释放而成为义子。\par

XXXII. \Person{毕达哥拉斯}认为，给事物命名的人应当被视为不仅最有智慧，而且是最年长的智者。因此，我们必须精确地查考圣经，因为它们被公认是用比喻表达的，并从名称中寻出圣神关于事物所提出、通过将祂的心意印在表达上（可以这样说）而教导的思想；这样，具有多种含义的名称，经过精确研究，得以阐明，那隐藏在多层覆盖之下的东西，经过处理和了解，就会显现出来并闪耀光辉。因为铅也是如此，摩擦时变白；白铅是由黑铅产生的。同样，\OriginalTerm{科学知识}{gnosis}将其光和明照耀在事物上，显明它自身就是真正的神圣智慧，纯净的光，照亮那些眼球清澈的人，使之确凿地看见和理解真理。\par

XXXIII. 于是，我们从那光的源头点燃火炬，凭借那以光为对象的炽热渴望，并尽力与之同化，我们就成了充满光的人，真正的以色列人。因为祂称那些以渴望和追求致力于相似神性的人为朋友和兄弟。\par

XXXIV. 洁净的地方和草地曾领受神圣异象的声音和异象。但每个完全净化的人，必被认为配得神圣教导和能力。\par

XXXV. 如今我知道，\OriginalTerm{知识}{gnosis}的奥迹在许多人看来是笑柄，尤其当它不用诡辩的比喻言语来修饰时。少数人起初对之惊讶，如同在黑暗的欢宴中突然亮光。随后，习惯于理性训练，仿佛因喜悦而欢欣鼓舞，他们\Emph{赞美}主。……因为快乐的本质是解脱痛苦，同样，知识的本质是除去无知。正如那些沉睡最深的人，因被极其清晰而固定的梦境所支配，自以为最清醒；同样，那些最无知的人自以为知道得最多。但那些从这沉睡和迷乱中醒来，举目朝向光明与真理的人是有福的。\par

XXXVI. 因此，对于想要有一个温顺的、将信德与渴望融合的门徒的人来说，同样有必要自己操练并不断独自研究，探究其思辨的真理；当他自以为正确时，要降下来面对关于邻近事物的问题。因为幼鸟在巢中尝试飞行，操练它们的翅膀。\par

XXXVII. 因为诺斯替的德性处处使人成为良善、温和、无害、无痛、有福，并准备好以最佳的方式与一切神圣者交往，以最佳的方式与人交往，同时既沉思又行动的神圣形象，并藉着爱使他成为善的爱慕者。因为善，在那里被智慧沉思和理解，在这里则通过信德藉着自制和正义付诸实行：在肉身中实行天使般的服务；在身体内圣化灵魂，如同在一个洁净无瑕的地方。\par

XXXVIII. 反对\Person{塔提安}，他说\BibleQuote{有光}是恳求的。\BibleRef{创 1:3}那么，如果祂是在恳求至高天主，祂又怎么说\BibleQuote{我是天主，除我以外没有别的神？}\BibleRef{依 44:6}我们已经说过，对亵渎、胡言、放肆的言辞是有惩罚的；这些被理性惩罚和管教。\par

XXXIX. 他还说，因着她们的头发和装饰，妇女被管理这些事务的权能惩罚；这权能也曾赐给\Person{参孙}头发中的力量；这权能惩罚那些藉装饰头发引诱人犯奸淫的妇女。\par

XL. 正如因善的流出，人成为良善；同样，他们成为邪恶。善是天主的审判，是区分信者与不信者，是预先的审判，以免陷入更大的审判——这审判是矫正。\par

XLI. 圣经说，被遗弃的婴儿被交付给一位护守天使，由他训练和养育。经文说：\Quote{他们将如这世上百岁的忠信者。}因此\Person{伯多禄}在默示录中说：\Quote{一道火焰从那些婴儿中跳出，击打妇女的眼睛。}因为义人像芦苇中的火花闪耀，将审判万民。\BibleRef{智 3:7}\par

XLII. \Quote{与圣者一起，你将成为圣。}\Quote{按你的赞颂，你的名被荣耀。}天主藉我们的知识和嗣业得荣耀。因此也说：\Quote{主活着，}以及\Quote{主复活了。}\BibleRef{路 24:34}\par

XLIII. \Quote{我所不认识的民族服事了我}——因盟约我不认识他们，外邦的子孙，他们贪恋属于别人的东西。\par

XLIV. \Quote{扩大祂君王的救恩。}所有忠信者都称为君王，藉嗣业被带进王权。\par

XLV. 忍耐是超越蜂蜜的甘甜；不是因为忍耐本身，而是因忍耐的果实。因此，自控的人没有激情，因为他克制激情，不是没有努力；但当习惯形成时，他不再是自控的人，因为他已受一种习惯和圣神的影响。\par

XLVI. 灵魂里的激情被称为神灵——不是权能的神灵，因为那样的话受激情影响的人会成为一群魔鬼；但它们\Emph{之所以如此称呼}，是因为它们所传递的冲动。因为灵魂本身，通过变化，采取这种或那种邪恶的品质，就被说成是接受神灵。\par

XLVII. 圣言并不命令我们放弃财产，而是管理财产而无过分的爱恋；无论发生什么，不烦恼不忧伤；不贪求获得。天主眷顾命令远离伴随激情的拥有，远离一切过分的爱恋，并\Emph{由此}使那些仍留在肉身中的人回转。\par

XLVIII. 例如，\Person{伯多禄}在默示录中说，流产的婴儿将享有更好的命运；他们被托付给一位护守天使，以致获知识后，能获得更好的居所，体验与若在身体里同样的经验。但其他的仅仅获得救恩，如同受伤害和怜悯，不受惩罚，领受这赏报。\par

XLIX. \Person{伯多禄}在默示录中说，妇女的乳汁从乳房流出并凝结，会产生微小的野兽，吞食肉体，再流回她们体内吞噬她们：教导说惩罚因罪恶而生。他说这些从罪恶产生；正如因他们的罪恶，人民被出卖。又因他们对基督缺乏信德，正如宗徒所说，他们被蛇咬。\par

一位古人曾说过，胚胎是有生命的；因为灵魂在子宫被洁净、为受孕准备之后进入子宫，并由一位掌管生育、知晓受孕时机、促使妇人交合的天使引入；并且，当种子被存下时，那在种子中的灵，可以说是被占有，因而在成形过程中被接纳结合。他引用这一点作为对所有人的证据：当天使向不育者报喜时，他们在受孕之前便引入灵魂。而在福音中\BibleQuote{胎儿跳跃}\BibleRef{路 1:43}，如同有生命之物。并且不育者之所以不育，是因为那为种子沉积而结合的灵魂没有被引入，以确保受孕和生育。\par

\BibleQuote{诸天陈述天主的荣耀。}诸天有多种含义，既指由空间和运转所界定的，也指由盟约所界定的——即首造天使的直接行动。因为盟约导致了天使更为特殊的显现——在亚当的情形，在诺厄的情形，在亚巴郎的情形，在梅瑟的情形。因为被主所推动，首造天使对那些依附于先知的天使施加影响，视盟约为天主的荣耀。此外，天使在地上所行之事，也是首造天使为了天主的荣耀而行的。\par

主首先被称为\Quote{诸天}，然后是首造者；在此之后，还有律法之前的圣人们，如族长、梅瑟和先知们；然后还有宗徒们。\BibleQuote{穹苍显示祂手的化工。}他用\Quote{穹苍}指称天主，那无情感且不改变的，正如同一达味在别处说：\BibleQuote{上主，我的力量，我的避难所，我爱慕你。}因此，穹苍本身显示祂手的化工——即显示并彰显祂天使的化工。因为祂显示并彰显了祂所造的那些。\par

\BibleQuote{日复一日发出言语。}既然诸天有多种含义，日也是如此。如今言语就是主；祂也常被称为日。\BibleQuote{夜复一夜显示知识。}魔鬼知道主要来。但他不相信祂是天主；因此他也试探祂，以便知道祂是否是大能的。据说，\BibleQuote{他离开祂，暂时离开了祂；}意思是，他把发现推迟到复活之后。因为他知道那将要复活的就是主。同样，邪魔们也是如此；因为连他们也曾怀疑撒罗满就是主，但在撒罗满犯罪时，他们就知道他不是。\BibleQuote{夜复一夜。}所有邪魔都知道那位在受难后复活的就是主。而且哈诺客已经说过，那些犯过罪的天使教给人类天文学和占卜，以及其余技艺。\par

\BibleQuote{没有言辞或话语，其声音不被听闻，}无论是日还是夜的。\BibleQuote{他们的声音传遍了全地。}祂将话语转至圣徒们身上，祂称他们为诸日和诸天。\par

星辰，即属灵的身体，与管辖它们的天使有交流，并受他们治理，它们并非万有产生的原因，而是正在发生、将要发生和已经发生之事的标记，涉及气候变迁、丰收与歉收、瘟疫与热病，以及人类之事。星辰丝毫不施加影响，只是指示什么是现在、将来和过去的事。\par

\BibleQuote{祂在太阳中安置了祂的帐幕。}这里有一种语序转换。因为所说的是第二次来临。因此，我们必须按正确的顺序阅读这倒装的语序：\BibleQuote{他如同新郎走出洞房，又如壮士欢然奔跑前路。从天的一端是他的出发；没有人能躲避他的炎热；}然后，\BibleQuote{祂在太阳中安置了祂的帐幕。}\par

有些人说，祂将主的身体安放在太阳中，如赫尔莫革内斯所言。而\Quote{祂的帐幕}，有些人说，是祂的身体，另一些人说是信众的教会。\par

我们的庞泰诺常说，预言表达时大多不定，用现在指未来，也用现在指过去。这也可以在此看到。因为\BibleQuote{祂安置了}既用于过去也用于未来。对于未来，是因为在这时期完成后（这时期将按目前的结构运行），主会来恢复义人，即信众（祂在他们内安息，如同在帐棚里）到同一个合一中；因为所有人都是一个身体，同一族类，并且选择了同一信德和义德。但有些人作为头，有些人作为眼，有些人作为耳，有些人作为手，有些人作为胸，有些人作为脚，将在太阳中被安置，发出光辉。\BibleQuote{犹如太阳发光}\BibleRef{玛 13:43}或\Quote{在太阳中}，因为有一位高位天使在太阳中。因为他被指派管辖白日，正如月亮管辖黑夜。\BibleRef{创 1:18}现在天使被称为日。据说，他们将与太阳中的天使一起被指派同一个居所，有时在某种程度上成为太阳，如同那单一身体的头。此外，他们也是日子的统治者，如同那太阳中的天使，为了更大的目的，他们在他之前迁移到同一地方。并且，他们将逐渐升迁，达到第一居所，与过去的\BibleQuote{祂安置了}相应：这样，首造天使不再按眷顾执行特定的服役，而是得以安息，专务默观天主；而次于他们的天使将升迁到他们留下的职位；再其下的天使也类似。\par

LVII. 那么，按照宗徒所说，有居于巅峰者，即最先受造者。他们就是宝座，虽是掌权者，作为最先受造者，因为天主在他们内安息，也如在那些相信的人内。因为各人按照自己进步的层次，以独特的方式拥有对天主的认识；天主在此认识中安息，那些拥有认识的人因认识而成为不朽的。难道不应当这样理解\BibleQuote{祂在太阳中搭起祂的帐幕}吗？天主\BibleQuote{安置在太阳中}，即安置在祂身旁的天主内，如同在福音中所说的：\BibleQuote{厄里、厄里，代替我的天主、我的天主}。而所谓\BibleQuote{超越一切统治、一切权柄、一切能力、一切可称呼的名号}的，就是从人类中得以成全为天使和总天使的人，以至于上升到最先受造之天使的本性。因为那些从人转变为天使的人，在达到成全之后，由天使教导一千年。然后，那些施教者被升到总天使的权柄；那些受教者又去教导那些从人转变为天使的人。这样，在后来的预定时期内，他们被带到与天使身体相应的状态。\par

LVIII. \BibleQuote{天主的法律是完善的，能转化灵魂。}救主本身被称为法律和圣言，如同伯多禄在\WorkTitle{宣讲}中所说，以及先知所说：\BibleQuote{因为法律将出自熙雍，上主的话将出自耶路撒冷。}\BibleRef{依 2:3}\par

LIX. \BibleQuote{上主的证言是可靠的，能使孩童有智慧。}主的盟约是真实的，使孩童有智慧；那些离弃邪恶的人，包括宗徒，然后也包括我们。此外，主的证言——即祂在苦难后复活，已由事实证实——引领教会坚定于信德。\par

LX. \BibleQuote{对上主的敬畏是纯洁的，永远常存。}他说那些从恐惧转向信德和正义的人永远常存。\par

\BibleQuote{上主的判断是真实的}——确实，且不可推翻；按照正义赐予赏报，带领义人达致信德的合一。因为这显现在\BibleQuote{因同一而成义}这话中。\BibleQuote{这样的愿望超越黄金和宝石。}\par

LXI. \BibleQuote{因为你的仆人也遵守它们。}这并非说只有达味被称为仆人；而是整个得救的百姓被称为天主的仆人，因着对诫命的服从。\par

LXII. \BibleQuote{求你洗净我隐秘的\Emph{过错}}——与正确理性相悖的思想——缺陷。因为祂称这为义人所不属之事。\par

LXIII. \BibleQuote{如果他们没有辖制我，那么我就纯真无过。}如果那像迫害主一样迫害我的人没有辖制我，我就不是纯真的。因为除非被人迫害，没有人成为殉道者；除非受委屈而不报复，无人显现为义；除非……忍耐……\par

\SourceRecord{Translated by William Wilson.；Ante-Nicene Fathers；Vol. 8.；Edited by Alexander Roberts, James Donaldson, and A. Cleveland Coxe.；Buffalo, NY: Christian Literature Publishing Co.,；1886.；<http://www.newadvent.org/fathers/0802.htm>.}
